% frontmatter/preface.tex
\chapter*{Preface}
\addcontentsline{toc}{chapter}{Preface}

Relational database design is a discipline of meaning.
It is about deciding what exists, how it relates, and what rules must always hold.
This book treats design as a conceptual activity before any SQL dialect or storage technology.
It assumes you want models that are stable, explainable, and correct.

The approach here is deliberately sober.
We use short sentences when clarity matters and longer ones when the idea needs context.
We use examples to anchor meaning, and constraints to make that meaning enforceable.
The goal is not to draw the prettiest diagrams.
The goal is to build models that survive reality.

This book is not about performance tuning, indexing, or query optimization.
Those are important, but they come later.
If the conceptual model is wrong, no amount of optimization will save it.
If the model is right, implementation becomes a disciplined translation.

The material grew from real projects in healthcare, education, and product systems.
The recurring lesson is that ambiguity is the real enemy.
A model that reads like a clear set of rules is already most of the way to a reliable system.

I hope this book helps you build that kind of model.
